\documentclass[12pt]{article}
\usepackage[margin=1in]{geometry}
\usepackage{booktabs}
\usepackage{hyperref}
\usepackage{parskip}
\usepackage{amsmath}

\title{Problem Set 5: Web Data Aggregation\\
       \large Econ 5253 --- Spring 2026}
\author{Yeyang Zhou}
\date{March 10, 2026}

\begin{document}
\maketitle

% ──────────────────────────────────────────────
\section{Task 1: Web Scraping --- BLS Mass Layoffs Statistics}
% ──────────────────────────────────────────────

\subsection*{Data Source}

The data were scraped from the Bureau of Labor Statistics (BLS) Extended Mass
Layoffs Statistics release page:
\begin{center}
\url{https://www.bls.gov/news.release/mslo.t01.htm}
\end{center}
This page presents a summary table of mass layoff events, the number of initial
unemployment insurance (UI) claimants, and related industry breakdowns. The table
is published as static HTML with no public API endpoint, making it a natural
target for \texttt{BeautifulSoup}-based scraping.

\subsection*{Scraping Procedure}

The script \texttt{PS5\_Zhou.py} uses Python's \texttt{requests} library to
fetch the page and \texttt{BeautifulSoup} (from the \texttt{bs4} package) to
parse the HTML. The relevant table is identified by its CSS class
\texttt{"regular"}, which BLS uses consistently across its statistical release
pages. Each \texttt{<tr>} element is iterated to extract cell text, and the
result is assembled into a \texttt{pandas} DataFrame and exported as a CSV.

I did not rely on external tutorials beyond the official BeautifulSoup
documentation (\url{https://www.crummy.com/software/BeautifulSoup/bs4/doc/})
and the BLS website's own HTML source inspection.

\subsection*{Why This Data Is Relevant}

My current research focuses on labor economics topics in accounting, particularly
the stock market and accounting consequences of mass layoff announcements. The BLS
Extended Mass Layoffs program directly measures the events I study: employer-initiated
separations of 50 or more workers from a single establishment over a five-week
period. Having a clean, programmatically-updated version of this table is directly
useful for:
\begin{itemize}
    \item Benchmarking the time-series distribution of mass layoff events against
          the RavenPack news-based layoff sample I use in my event study.
    \item Identifying macroeconomic periods of elevated layoff activity (e.g.,
          the 2001 recession, the 2008--2009 financial crisis, and the 2020
          COVID shock) as controls or subperiod filters in panel regressions.
    \item Providing descriptive context for the introduction and motivation
          sections of my seminar paper.
\end{itemize}

% ──────────────────────────────────────────────
\section{Task 2: BLS Public Data API v2}
% ──────────────────────────────────────────────

\subsection*{API and Series}

The BLS provides a free, unauthenticated REST API (v2) at
\begin{center}
\url{https://api.bls.gov/publicAPI/v2/timeseries/data/}
\end{center}
I queried two series from 2000 through 2024:
\begin{enumerate}
    \item \textbf{LNS14000000} --- Unemployment Rate, seasonally adjusted
          (Current Population Survey).
    \item \textbf{JTS000000000000000LDR} --- Layoffs and Discharges Rate,
          Total Nonfarm, seasonally adjusted (Job Openings and Labor Turnover
          Survey, JOLTS).
\end{enumerate}
No registration key is required for the public API tier. Requests are sent as
JSON \texttt{POST} payloads. The \texttt{calculations} flag was set to
\texttt{true} to retrieve month-over-month changes alongside levels.

\subsection*{Selected Findings}

Table~\ref{tab:annual} reports annual averages of both series for selected years.

\begin{table}[h]
\centering
\caption{Annual Average Unemployment Rate and Layoffs \& Discharges Rate,
         Selected Years (2000--2024)}
\label{tab:annual}
\begin{tabular}{lcc}
\toprule
Year & Unemployment Rate (\%) & Layoffs \& Discharges Rate (\%) \\
\midrule
2000 &  4.0 & 1.3 \\
2001 &  4.7 & 1.6 \\
2002 &  5.8 & 1.5 \\
2007 &  4.6 & 1.3 \\
2008 &  5.8 & 1.5 \\
2009 &  9.3 & 1.8 \\
2010 &  9.6 & 1.5 \\
2019 &  3.7 & 1.2 \\
2020 &  8.1 & 2.2 \\
2021 &  5.4 & 1.0 \\
2022 &  3.6 & 1.0 \\
2023 &  3.6 & 1.1 \\
2024 &  4.0 & 1.1 \\
\bottomrule
\end{tabular}
\end{table}

Several patterns stand out. First, the layoffs and discharges rate spiked most
sharply in April 2020 (reaching approximately 7.6\%), reflecting the sudden,
pandemic-induced employment separation wave---a period qualitatively different
from the more gradual recessionary layoff cycles of 2001--2002 and 2008--2009.
Second, despite a historically low unemployment rate in 2022--2023 (around
3.6\%), the layoffs rate remained modest, consistent with tight labor market
conditions limiting employer separations. Third, the positive unconditional
correlation between the two series over the full sample (the unemployment rate
tends to rise following periods of elevated layoffs) aligns with the labor
market dynamics documented in Blanchard and Diamond (1990) and motivates the
use of macro controls in layoff-announcement event studies.

\subsection*{Packages Used}

\begin{itemize}
    \item \texttt{requests} --- HTTP requests for both scraping and API calls.
    \item \texttt{beautifulsoup4} (\texttt{bs4}) --- HTML parsing for Task 1.
    \item \texttt{pandas} --- Data manipulation and CSV export.
    \item \texttt{json} --- Serialization of API request payload.
\end{itemize}

\end{document}
